\SourcePage{330}{335}
\section{The electron propagator}

The propagation functions, or propagators, introduced in the preceding sections play a central role in the apparatus of quantum electrodynamics. The photon propagator $D_{\mu\nu}$ becomes the principal quantity characterizing the interaction of two electrons. This role is clearly displayed by its place in the electron-scattering amplitude, where $D_{\mu\nu}$ is multiplied by the transition currents of the two particles. A simi\SourcePageJoin{331}{336}lar role is played by the electron propagator in the interaction between an electron and a photon.

We now turn to the actual calculation of the propagators, beginning with the electron case.

Apply the operator $\gamma\hat p-m$, where $\hat p_\mu=i\partial_\mu$, to the function
\begin{equation}
G_{ik}(x-x')=-i\langle0|\mathop{\mathrm T}\hat\psi_i(x)\widehat{\bar\psi}_k(x')|0\rangle
\tag{75.1}
\label{eq:75.1}
\end{equation}
($i$ and $k$ are bispinor indices). Since $\hat\psi(x)$ satisfies the Dirac equation $(\gamma\hat p-m)\hat\psi(x)=0$, the result is zero at every $x$ except where $t=t'$. Indeed, $G(x-x')$ approaches different limits as $t\to t'+0$ and $t\to t'-0$. By the definition (74.8), these limits are respectively
\[
-i\langle0|\hat\psi_i(\mathbf r,t)\widehat{\bar\psi}_k(\mathbf r',t)|0\rangle
\quad\text{and}\quad
+i\langle0|\widehat{\bar\psi}_k(\mathbf r',t)\hat\psi_i(\mathbf r,t)|0\rangle.
\]
As we shall see, they do not coincide on the light cone. The derivative $\partial G/\partial t$ therefore acquires an additional delta-function term:
\begin{equation}
\begin{aligned}
\frac{\partial G}{\partial t}={}&-i\left\langle0\left|\mathop{\mathrm T}
\frac{\partial\hat\psi_i(x)}{\partial t}\widehat{\bar\psi}_k(x')\right|0\right\rangle\\
&+\delta(t-t')\bigl(G|_{t\to t'+0}-G|_{t\to t'-0}\bigr).
\end{aligned}
\tag{75.2}
\end{equation}
Noting that the time derivative enters $\gamma\hat p-m$ as $i\gamma^0\partial/\partial t$, we thus have
\begin{equation}
(\gamma\hat p-m)_{ik}G_{kl}(x-x')
=\delta(t-t')\gamma^0_{ik}\langle0|\{\hat\psi_k(\mathbf r,t),
\widehat{\bar\psi}_l(\mathbf r',t)\}_+|0\rangle.
\tag{75.3}
\end{equation}

Let us calculate the anticommutator in this expression. Multiplying the operators $\hat\psi(\mathbf r,t)$ and $\widehat{\bar\psi}(\mathbf r',t)$ (see (73.6)) and applying the interchange rules for the fermion operators $\hat a_{\mathbf p}$ and $\hat b_{\mathbf p}$, we find
\begin{equation}
\{\hat\psi_i(\mathbf r,t),\widehat{\bar\psi}_k(\mathbf r',t)\}_+
=\sum_{\mathbf p}\bigl[\psi_{\mathbf p i}(\mathbf r)\bar\psi_{\mathbf p k}(\mathbf r')
+\psi_{-\mathbf p i}(\mathbf r)\bar\psi_{-\mathbf p k}(\mathbf r')\bigr].
\tag{75.4}
\end{equation}
Here $\psi_{\pm\mathbf p}(\mathbf r)$ are the wave functions without the time factor (as in \S\S~73 and 74, their polarization indices are suppressed for brevity). The collection of all $\psi_{\pm\mathbf p}(\mathbf r)$, the eigenfunctions of the electron Hamiltonian, forms a complete normalized system; hence, by the general properties of such systems (compare III, (5.12)),
\begin{equation}
\sum_{\mathbf p}\bigl[\psi_{\mathbf p i}(\mathbf r)\psi^*_{\mathbf p k}(\mathbf r')
+\psi_{-\mathbf p i}(\mathbf r)\psi^*_{-\mathbf p k}(\mathbf r')\bigr]
=\delta_{ik}\delta(\mathbf r-\mathbf r').
\tag{75.5}
\end{equation}

\SourcePage{332}{337}
The sum on the right-hand side of (75.4) differs from this one by the replacement of $\psi_k^*$ with $(\psi^*\gamma^0)_k$ and is therefore equal to $\gamma^0_{ik}\delta(\mathbf r-\mathbf r')$. Thus
\begin{equation}
\{\hat\psi_i(\mathbf r,t),\widehat{\bar\psi}_k(\mathbf r',t)\}_+
=\delta(\mathbf r-\mathbf r')\gamma^0_{ik}.
\tag{75.6}
\end{equation}
In particular, this formula implies the statement already mentioned in \S~74 that $\hat\psi$ and $\widehat{\bar\psi}$ anticommute outside the light cone. If $(x-x')^2<0$, a reference frame always exists in which $t=t'$; when $\mathbf r\ne\mathbf r'$ as well, the anticommutator (75.6) is indeed zero.

Substitution of (75.6) into (75.3), with the bispinor indices suppressed, finally gives\footnote{Written explicitly with bispinor indices,
\[
(\gamma\hat p-m)_{il}G_{lk}(x-x')=\delta^{(4)}(x-x')\delta_{ik}.
\tag{75.7a}
\]}
\begin{equation}
(\gamma\hat p-m)G(x-x')=\delta^{(4)}(x-x').
\tag{75.7}
\end{equation}
Thus the electron propagator satisfies the Dirac equation with a delta function on the right-hand side. In other words, it is a \emph{Green function} for the Dirac equation.

In what follows we shall need not $G(\xi)$ itself, where $\xi=x-x'$, but its Fourier components
\begin{equation}
G(p)=\int G(\xi)e^{ip\xi}\,d^4\xi.
\tag{75.8}
\end{equation}
(the propagator in the momentum representation). Taking the Fourier component of both sides of (75.7), we find that $G(p)$ satisfies the system of algebraic equations
\begin{equation}
(\gamma p-m)G(p)=1.
\tag{75.9}
\end{equation}
The solution of this system is
\begin{equation}
G(p)=\frac{\gamma p+m}{p^2-m^2}.
\tag{75.10}
\end{equation}
The four components of the 4-vector $p$ in $G(p)$ are independent variables; they are not constrained by $p^2=p_0^2-\mathbf p^2=m^2$. Writing the denominator of (75.10) as $p_0^2-(\mathbf p^2+m^2)$ shows that, as a function of $p_0$ at fixed $\mathbf p$, $G(p)$ has two poles: $p_0=\pm\varepsilon$, where $\varepsilon=\sqrt{\mathbf p^2+m^2}$. Integration over $dp_0$ in the integral
\begin{equation}
\begin{aligned}
G(\xi)&=\frac{1}{(2\pi)^4}\int e^{-ip\xi}G(p)\,d^4p\\
&=\frac{1}{(2\pi)^4}\int d^3p\,e^{i\mathbf p\boldsymbol\xi}
\int dp_0\,e^{-ip_0\tau}G(p)
\end{aligned}
\tag{75.11}
\end{equation}
\SourcePage{333}{338}
where $\tau=t-t'$, consequently raises the question of how the poles are to be bypassed; until this prescription is specified, (75.10) is still essentially undefined.

To settle this question, return to the original definition (75.1). Substitute the expansions (73.6) for the $\psi$-operators, observing that only the following vacuum mean values of products of creation and annihilation operators are nonzero:
\[
\langle0|\hat a_{\mathbf p}\hat a_{\mathbf p}^+|0\rangle=1,
\qquad
\langle0|\hat b_{\mathbf p}\hat b_{\mathbf p}^+|0\rangle=1.
\]
(Because the vacuum contains no particles, a particle must first be created by $\hat a_{\mathbf p}^+$ or $\hat b_{\mathbf p}^+$ before it can be annihilated by $\hat a_{\mathbf p}$ or $\hat b_{\mathbf p}$.) We obtain
\begin{equation}
\begin{aligned}
G_{ik}(x-x')&=-i\sum_{\mathbf p}\psi_{\mathbf p i}(\mathbf r,t)
\bar\psi_{\mathbf p k}(\mathbf r',t')\\
&=-i\sum_{\mathbf p}e^{-i\varepsilon(t-t')}
\psi_{\mathbf p i}(\mathbf r)\bar\psi_{\mathbf p k}(\mathbf r'),
\qquad t-t'>0;
\end{aligned}
\tag{75.12}
\end{equation}
\[
\begin{aligned}
G_{ik}(x-x')&=i\sum_{\mathbf p}\bar\psi_{-\mathbf p k}(\mathbf r',t')
\psi_{-\mathbf p i}(\mathbf r,t)\\
&=i\sum_{\mathbf p}e^{i\varepsilon(t-t')}
\psi_{-\mathbf p i}(\mathbf r)\bar\psi_{-\mathbf p k}(\mathbf r'),
\qquad t-t'<0.
\end{aligned}
\]
(For $t>t'$ only the electron terms contribute to $G$, and for $t<t'$ only the positron terms.)

Replace the sum over $\mathbf p$ conceptually by an integral over $d^3p$ and compare (75.12) with (75.11). The integral
\begin{equation}
\int e^{-ip_0\tau}G(p)\,dp_0
\tag{75.13}
\end{equation}
must have the phase factor $e^{-i\varepsilon\tau}$ for $\tau>0$ and $e^{i\varepsilon\tau}$ for $\tau<0$. This condition is met if the poles $p_0=\varepsilon$ and $p_0=-\varepsilon$ are bypassed above and below, respectively, in the complex $p_0$ plane:
\begin{equation}
\begin{tikzpicture}[baseline=-.45ex,x=.8cm,y=.8cm,>=Stealth,line width=.55pt,font=\small]
\draw[->] (-3.25,0)--(-1.25,0);
\draw (-1.25,0) arc[start angle=180,end angle=360,radius=.45];
\draw[->] (-.35,0)--(1.15,0);
\draw (1.15,0) arc[start angle=180,end angle=0,radius=.45];
\draw[->] (2.05,0)--(3.25,0);
\fill (-.8,0) circle (1.2pt);
\fill (0,0) circle (1.2pt);
\fill (1.6,0) circle (1.2pt);
\node[above=3pt] at (-.8,0) {$-\varepsilon$};
\node[below=3pt] at (0,0) {$0$};
\node[below=3pt] at (1.6,0) {$+\varepsilon$};
\end{tikzpicture}
\tag{75.14}
\end{equation}
Indeed, for $\tau>0$ the integration path is closed by a semicircle at infinity in the lower half-plane, so the value of (75.13) is given by the residue at the pole $p_0=+\varepsilon$. For $\tau<0$ the contour is closed in the upper half-plane and the integral is determined by the residue at $p_0=-\varepsilon$. The required result follows in both cases.

\SourcePage{334}{339}
This bypass prescription, the \emph{Feynman prescription}, can be stated in another way: integration is performed everywhere along the real axis itself, but the particle mass $m$ is assigned an infinitesimal negative imaginary part:
\begin{equation}
m\to m-i0.
\tag{75.15}
\end{equation}
Indeed, then
\[
\varepsilon\to\sqrt{\mathbf p^2+(m-i0)^2}
=\sqrt{\mathbf p^2+m^2-i0}=\varepsilon-i0.
\]
In other words, the poles $p_0=\pm\varepsilon$ are displaced below and above the real axis:
\begin{equation}
\begin{tikzpicture}[baseline=-.45ex,x=.8cm,y=.8cm,>=Stealth,line width=.55pt,font=\small]
\draw[->] (-3.2,0)--(-.45,0);
\draw[->] (-.45,0)--(3.2,0);
\fill (-1.35,.42) circle (1.4pt);
\fill (0,0) circle (1.4pt);
\fill (1.35,-.42) circle (1.4pt);
\node[above=3pt] at (-1.35,.42) {$-\varepsilon+i0$};
\node[below=3pt] at (0,0) {$0$};
\node[below left=2pt] at (1.35,-.42) {$+\varepsilon$};
\node[below right=2pt] at (1.35,-.42) {$-i0$};
\end{tikzpicture}
\tag{75.16}
\end{equation}
so that integration along this axis becomes equivalent to integration along the path in (75.14)\footnote{It is useful to note that the pole-displacement prescription corresponds to an infinitesimal damping of $G(x-x')$ with $|\tau|$, where $\tau=t-t'$. Indeed, if the values of $p_0$ at the displaced poles are written as $-(\varepsilon-i\delta)$ and $+(\varepsilon-i\delta)$, where $\delta\to+0$, the time factor in (75.13) is $\exp(-i\varepsilon|\tau|-\delta|\tau|)$.}. With the prescription (75.15), the propagator (75.10) may be written
\begin{equation}
G(p)=\frac{\gamma p+m}{p^2-m^2+i0}.
\tag{75.17}
\end{equation}
The integration rule for a displaced pole is expressed by
\begin{equation}
\frac{1}{x+i0}=\mathop{\mathrm P}\frac{1}{x}-i\pi\delta(x).
\tag{75.18}
\end{equation}
This is to be understood in the sense that multiplication by a function $f(x)$ and integration gives
\begin{equation}
\int_{-\infty}^{\infty}\frac{f(x)}{x+i0}\,dx
=\mathop{\mathrm P}\int_{-\infty}^{\infty}\frac{f(x)}{x}\,dx-i\pi f(0),
\tag{75.19}
\end{equation}
where the crossed integral sign, or the symbol $\mathrm P$, denotes the principal value.

The Green function (75.10) is the product of the bispinor factor $\gamma p+m$ and the scalar
\begin{equation}
G^{(0)}(p)=\frac{1}{p^2-m^2}.
\tag{75.20}
\end{equation}
The corresponding coordinate function $G^{(0)}(\xi)$ is evidently a solution of
\begin{equation}
(\hat p^2-m^2)G^{(0)}(x-x')=\delta^{(4)}(x-x'),
\tag{75.21}
\end{equation}
\SourcePage{335}{340}
that is, a Green function for the equation $(\hat p^2-m^2)\psi=0$. In this sense, $G^{(0)}(x-x')$ may be called the propagator for scalar particles. A direct calculation similar to the one above readily shows that the scalar-field propagation function is expressed in terms of the $\psi$-operators (11.2) by
\begin{equation}
G^{(0)}(x-x')=-i\langle0|\mathop{\mathrm T}\hat\psi(x)\hat\psi^+(x')|0\rangle,
\tag{75.22}
\end{equation}
in analogy with the definition (75.1). Here the chronological product is defined, as for all boson operators, by
\begin{equation}
\mathop{\mathrm T}\hat\psi(x)\hat\psi^+(x')=
\begin{cases}
\hat\psi(x)\hat\psi^+(x'),&t>t',\\
\hat\psi^+(x')\hat\psi(x),&t<t'
\end{cases}
\tag{75.23}
\end{equation}
with the same signs for $t>t'$ and $t<t'$.
